\section{Injective modules}
\label{section-injectives-modules}

\noindent
Some lemmas on injective modules.

\begin{definition}
\label{definition-projective}
Let $R$ be a ring. An $R$-module $J$ is {\it injective} if and only if
the functor $\Hom_R(-, J) : \text{Mod}_R \to \text{Mod}_R$ is
an exact functor.
\end{definition}

\noindent
The functor $\Hom_R(- , M)$ is left exact for any $R$-module $M$, see
Algebra, Lemma \ref{algebra-lemma-hom-exact}.
Hence the condition for $J$ to be injective really signifies that given
an injection of $R$-modules $M \to M'$ the map
$\Hom_R(M', J) \to \Hom_R(M, J)$ is surjective.

\medskip\noindent
Before we reformulate this in terms of ${Ext}$-modules we discuss the
relationship between $\Ext^1_R(M, N)$ and extensions as in
Homology, Section \ref{homology-section-extensions}.

\begin{lemma}
\label{lemma-relation-ext-ext}
Let $R$ be a ring. Let $\mathcal{A}$ be the abelian category of
$R$-modules. There is a canonical isomorphism
$\Ext_\mathcal{A}(M, N) = \Ext^1_R(M, N)$
compatible with the long exact sequences of
Algebra, Lemmas \ref{algebra-lemma-long-exact-seq-ext} and
\ref{algebra-lemma-reverse-long-exact-seq-ext}
and the $6$-term exact sequences of
Homology, Lemma \ref{homology-lemma-six-term-sequence-ext}.
\end{lemma}

\begin{proof}
Omitted.
\end{proof}

\begin{lemma}
\label{lemma-characterize-injective}
Let $R$ be a ring. Let $J$ be an $R$-module.
The following are equivalent
\begin{enumerate}
\item $J$ is injective,
\item $\Ext^1_R(M, J) = 0$ for every $R$-module $M$.
\end{enumerate}
\end{lemma}

\begin{proof}
Let $0 \to M'' \to M' \to M \to 0$ be a short exact sequence of $R$-modules.
Consider the long exact sequence
$$
\begin{matrix}
0
\to \Hom_R(M, J)
\to \Hom_R(M', J)
\to \Hom_R(M'', J)
\\
\phantom{0\ }
\to \Ext^1_R(M, J)
\to \Ext^1_R(M', J)
\to \Ext^1_R(M'', J)
\to \ldots
\end{matrix}
$$
of Algebra, Lemma \ref{algebra-lemma-reverse-long-exact-seq-ext}.
Thus we see that (2) implies (1). Conversely, if $J$ is injective
then the $\Ext$-group is zero by
Homology, Lemma \ref{homology-lemma-characterize-injectives} and
Lemma \ref{lemma-relation-ext-ext}.
\end{proof}

\begin{lemma}
\label{lemma-characterize-injective-bis}
Let $R$ be a ring. Let $J$ be an $R$-module.
The following are equivalent
\begin{enumerate}
\item $J$ is injective,
\item $\Ext^1_R(R/I, J) = 0$ for every ideal $I \subset R$, and
\item for an ideal $I \subset R$ and module map $I \to J$
there exists an extension $R \to J$.
\end{enumerate}
\end{lemma}

\begin{proof}
If $I \subset R$ is an ideal, then the short exact sequence
$0 \to I \to R \to R/I \to 0$ gives an exact sequence
$$
\Hom_R(R, J) \to
\Hom_R(I, J) \to
\Ext^1_R(R/I, J) \to 0
$$
by Algebra, Lemma \ref{algebra-lemma-reverse-long-exact-seq-ext}
and the fact that $\Ext^1_R(R, J) = 0$ as $R$ is projective
(Algebra, Lemma \ref{algebra-lemma-characterize-projective}).
Thus (2) and (3) are equivalent. In this proof
we will show that (1) $\Leftrightarrow$ (3) which is known
as Baer's criterion.

\medskip\noindent
Assume (1). Given a module map $I \to J$ as in (3) we find
the extension $R \to J$ because the map
$\Hom_R(R, J) \to \Hom_R(I, J)$ is surjective
by definition.

\medskip\noindent
Assume (3). Let $M \subset N$ be an inclusion of $R$-modules.
Let $\varphi : M \to J$ be a homomorphism. We will show that $\varphi$
extends to $N$ which finishes the proof of the lemma.
Consider the set of homomorphisms $\varphi' : M' \to J$
with $M \subset M' \subset N$ and $\varphi'|_M = \varphi$.
Define $(M', \varphi') \geq (M'', \varphi'')$ if
and only if $M' \supset M''$ and $\varphi'|_{M''} = \varphi''$.
If $(M_i, \varphi_i)_{i \in I}$ is a totally
ordered collection of such pairs, then we obtain a map
$\bigcup_{i \in I} M_i \to J$ defined by $a \in M_i$
maps to $\varphi_i(a)$. Thus Zorn's lemma applies.
To conclude we have to show that if the pair
$(M', \varphi')$ is maximal then $M' = N$.
In other words, it suffices to show, given
any subgroup $M \subset N$, $M \not = N$ and
any $\varphi : M \to J$, then we can find
$\varphi' : M' \to J$ with $M \subset M' \subset N$
such that (a) the inclusion $M \subset M'$ is strict, and
(b) the morphism $\varphi'$ extends $\varphi$.

\medskip\noindent
To prove this, pick $x \in N$, $x \not \in M$.
Let $I = \{f \in R \mid fx \in M\}$. This is an ideal of $R$.
Define a homomorphism $\psi : I \to J$ by $f \mapsto \varphi(fx)$.
Extend to a map $\tilde\psi : R \to J$ which is possible by assumption (3).
By our choice of $I$ the kernel of
$M \oplus R \to J$, $(y, f) \mapsto \varphi(y) + \tilde\psi(f)$
contains the kernel of the map $M \oplus R \to N$,
$(y, f) \mapsto y + fx$. Hence this homomorphism factors
through the image $M' = M + Rx$ and this extends the given homomorphism
as desired.
\end{proof}

\noindent
In the rest of this section we prove that there are enough injective
modules over a ring $R$. We start with the fact that $\mathbf{Q}/\mathbf{Z}$
is an injective abelian group. This follows from
Lemma \ref{lemma-injective-abelian}.

\begin{definition}
\label{definition-simple-functors}
Let $R$ be a ring.
\begin{enumerate}
\item For any $R$-module $M$ over $R$ we denote
$M^\vee = \Hom(M, \mathbf{Q}/\mathbf{Z})$
with its natural $R$-module structure. We think
of {\it $M \mapsto M^\vee$} as a contravariant functor
from the category of $R$-modules to itself.
\item For any $R$-module $M$ we denote
$$
F(M) = \bigoplus\nolimits_{m \in M} R[m]
$$
the {\it free module} with basis given by the elements $[m]$ with
$m \in M$. We let $F(M)\to M$, $\sum f_i [m_i] \mapsto \sum f_i m_i$
be the natural surjection of $R$-modules.
We think of $M \mapsto (F(M) \to M)$ as a functor from
the category of $R$-modules to the category of
arrows in $R$-modules.
\end{enumerate}
\end{definition}

\begin{lemma}
\label{lemma-vee-exact}
Let $R$ be a ring.
The functor $M \mapsto M^\vee$ is exact.
\end{lemma}

\begin{proof}
This because $\mathbf{Q}/\mathbf{Z}$
is an injective abelian group by Lemma \ref{lemma-injective-abelian}.
\end{proof}

\noindent
There is a canonical map $ev : M \to (M^\vee)^\vee$
given by evaluation: given $x \in M$ we let
$ev(x) \in (M^\vee)^\vee = \Hom(M^\vee, \mathbf{Q}/\mathbf{Z})$
be the map $\varphi \mapsto \varphi(x)$.

\begin{lemma}
\label{lemma-ev-injective}
For any $R$-module $M$ the evaluation map
$ev : M \to (M^\vee)^\vee$ is injective.
\end{lemma}

\begin{proof}
You can check this using that $\mathbf{Q}/\mathbf{Z}$ is an injective
abelian group. Namely, if $x \in M$ is not zero, then let
$M' \subset M$ be the cyclic group it generates. There exists
a nonzero map $M' \to \mathbf{Q}/\mathbf{Z}$ which necessarily does
not annihilate $x$. This extends to
a map $\varphi : M \to \mathbf{Q}/\mathbf{Z}$
and then $ev(x)(\varphi) = \varphi(x) \not = 0$.
\end{proof}

\noindent
The canonical surjection $F(M) \to M$ of $R$-modules turns into
a canonical injection, see above, of $R$-modules
$$
(M^\vee)^\vee \longrightarrow (F(M^\vee))^\vee.
$$
Set $J(M) = (F(M^\vee))^\vee$. The composition of $ev$ with this
the displayed map gives $M \to J(M)$ functorially in $M$.

\begin{lemma}
\label{lemma-JM-injective}
Let $R$ be a ring. For every $R$-module $M$ the
$R$-module $J(M)$ is injective.
\end{lemma}

\begin{proof}
Note that $J(M) \cong \prod_{\varphi \in M^\vee} R^\vee$ as an $R$-module.
As the product of injective modules is injective, it suffices to
show that $R^\vee$ is injective. For this we use that
$$
\Hom_R(N, R^\vee) =
\Hom_R(N, \Hom_{\mathbf{Z}}(R, \mathbf{Q}/\mathbf{Z})) =
N^\vee
$$
and the
fact that $(-)^\vee$ is an exact functor by Lemma
\ref{lemma-vee-exact}.
\end{proof}

\begin{lemma}
\label{lemma-injectives-modules}
Let $R$ be a ring.
The construction above defines a covariant functor
$M \mapsto (M \to J(M))$ from the category of
$R$-modules to the category of arrows of $R$-modules
such that for every module $M$ the output
$M \to J(M)$ is an injective map of $M$ into
an injective $R$-module $J(M)$.
\end{lemma}

\begin{proof}
Follows from the above.
\end{proof}

\noindent
In particular, for any map of $R$-modules $M \to N$
there is an associated morphism $J(M) \to J(N)$
making the following diagram commute:
$$
\xymatrix{
M \ar[d] \ar[r] & N \ar[d] \\
J(M) \ar[r] & J(N) }
$$
This is the kind of construction we would like to have in general.
In Homology, Section \ref{homology-section-injectives}
we introduced terminology to express this. Namely,
we say this means that the category of $R$-modules
has functorial injective embeddings.











